\documentclass[11pt]{article}
\usepackage[final]{acl}
\usepackage{times}
\usepackage{latexsym}
\usepackage{fontspec}
\usepackage{polyglossia}
\setdefaultlanguage{russian}
\setotherlanguage{english}
\setmainfont{Times New Roman}
\newfontfamily\cyrillicfont{Times New Roman}
\usepackage{microtype}
\usepackage{inconsolata}
\usepackage{graphicx}

\title{Формализация критериев оценки веб-поиска и сравнение современных систем}

\author{Игорь Пивнев, ИТМО}

\begin{document}
\maketitle

\begin{abstract}
В работе рассматриваются критерии оценки качества веб-поисковых систем и проводится
сравнительный анализ традиционных поисковиков (Google, Yandex) и инструментов на
базе больших языковых моделей (ChatGPT, LLM в поиске). Предлагается набор критериев
и сценариев использования, а также анализируются сильные и слабые стороны различных
подходов к поиску информации.
\end{abstract}

\section{Постановка задачи}

Целью данной работы является формализация критериев оценки веб-поисковых систем
и сравнение классического поиска с LLM-подходом. Особый интерес представляет различие
в поведении систем в зависимости от типа запроса и пользовательского сценария.

\section{Описание данных}

Анализ проводится на основе результатов поисковой выдачи для трех типов запросов:

\begin{itemize}
    \item фактологический: \textit{"когда вышел фильм аватар 3"},
    \item сравнительный: \textit{"сходить в кино в спб"},
    \item исследовательский: \textit{"что такое RAG"}.
\end{itemize}

Все запросы выполнялись на русском языке в рамках одного региона, что позволяет
наблюдать влияние региональной специфики на выдачу.

\section{Методы и инструменты}

\subsection{Критерии оценки}

В работе используются следующие критерии:

\begin{itemize}
    \item \textbf{Скорость выдачи} — субъективно оценивается как комфортная, если
    ответ появляется менее чем за секунду.
    
    \item \textbf{Качество поиска} — определяется наличием ответа на запрос среди
    первых пяти результатов выдачи.
    
    \item \textbf{Качество источников} — оценивается по их надежности и удобству:
    предпочтение отдается известным ресурсам (например, энциклопедиям или
    популярным платформам) и сайтам с понятной структурой.
    
    \item \textbf{Интерфейс} — важна простота и отсутствие перегруженности.
    
    \item \textbf{Помощь в формулировке запроса} — наличие автодополнения и подсказок.
    
    \item \textbf{Представление результатов} — наличие краткого ответа и удобной
    структуры выдачи.
\end{itemize}

\subsection{Сценарии поиска и пользователи}

Выделяются три сценария поиска:

\begin{itemize}
    \item фактологический — быстрый доступ к точному ответу,
    \item сравнительный — выбор между альтернативами,
    \item исследовательский — глубокое изучение темы.
\end{itemize}

Требования к поиску различаются в зависимости от пользователя. Новичкам важны
скорость и простота интерфейса, тогда как более опытные пользователи ожидают
разнообразия источников и глубины материалов.

\section{Результаты и интерпретация}

\subsection{Сравнение Google и Yandex}

Google демонстрирует более стабильные результаты в фактологическом поиске:
ответ на запрос чаще находится среди первых ссылок и представлен через
надежные источники.

Yandex сильнее проявляет себя в региональных и повседневных сценариях.
Например, для запроса о досуге он активно предлагает сервисы (такие как покупка
билетов), что делает поиск более прикладным. В этом контексте подобное поведение
скорее является преимуществом.

Интерфейс Google остается более минималистичным, тогда как Yandex предлагает
большее количество визуальных элементов и сервисов, что может восприниматься
как перегруженность.

Также наблюдается различие в характере источников: Google чаще показывает
международные и более формальные материалы, тогда как Yandex — более
локальные и разнообразные источники из СНГ.

\subsection{Сравнение с LLM-инструментами}

Встроенные LLM-решения в поисковых системах предоставляют краткие ответы
на основе найденных источников. Их качество в целом сопоставимо, однако
дополнительные функции (например, подсказки для следующих запросов в Yandex)
помогают пользователю продолжать поиск.

ChatGPT, в отличие от поисковых систем, лучше подходит для исследовательских
задач. Например, запрос \textit{"что такое RAG"} с последующими уточнениями
позволяет глубже разобраться в теме без перехода между источниками.

Однако в задачах, требующих актуальной информации (например, новости),
классические поисковики показывают более релевантные результаты.

\section{Анализ и обсуждение}

Классические поисковые системы остаются наиболее эффективными в задачах,
где важна актуальность информации и доступ к разнообразным источникам.

LLM-инструменты, в свою очередь, упрощают исследовательский процесс за счет
диалогового взаимодействия и способности агрегировать знания.

Таким образом, эти подходы не конкурируют напрямую, а дополняют друг друга.
Наиболее перспективным направлением является их интеграция, при которой
LLM выступает как слой поверх традиционного поиска.

На практике оптимальным решением является использование гибридного подхода,
например, поисковой системы с встроенными LLM-возможностями.

\end{document}
